% Options for packages loaded elsewhere
\PassOptionsToPackage{unicode}{hyperref}
\PassOptionsToPackage{hyphens}{url}
\PassOptionsToPackage{dvipsnames,svgnames,x11names}{xcolor}
\documentclass[
  11pt,
  a4paper,
]{article}
\usepackage{xcolor}
\usepackage[margin=1in]{geometry}
\usepackage{amsmath,amssymb}
\setcounter{secnumdepth}{-\maxdimen} % remove section numbering
\usepackage{iftex}
\ifPDFTeX
  \usepackage[T1]{fontenc}
  \usepackage[utf8]{inputenc}
  \usepackage{textcomp} % provide euro and other symbols
\else % if luatex or xetex
  \usepackage{unicode-math} % this also loads fontspec
  \defaultfontfeatures{Scale=MatchLowercase}
  \defaultfontfeatures[\rmfamily]{Ligatures=TeX,Scale=1}
\fi
\usepackage{lmodern}
\ifPDFTeX\else
  % xetex/luatex font selection
  \ifLuaTeX
    \usepackage{luaotfload}
    \directlua{luaotfload.add_fallback("mainfontfallback",{
      "NotoSerifBengali:mode=harf"
    })}
  \fi
  \setmainfont[,RawFeature={fallback=mainfontfallback}]{Latin Modern Roman}
  \ifLuaTeX
    \usepackage{luaotfload}
    \directlua{luaotfload.add_fallback("sansfontfallback",{
      "NotoSansBengali:mode=harf"
    })}
  \fi
  \setsansfont[,RawFeature={fallback=sansfontfallback}]{Latin Modern Sans}
  \ifLuaTeX
    \usepackage{luaotfload}
    \directlua{luaotfload.add_fallback("monofontfallback",{
      "NotoSansBengali:mode=harf"
    })}
  \fi
  \setmonofont[,RawFeature={fallback=monofontfallback}]{DejaVu Sans Mono}
\fi
% Use upquote if available, for straight quotes in verbatim environments
\IfFileExists{upquote.sty}{\usepackage{upquote}}{}
\IfFileExists{microtype.sty}{% use microtype if available
  \usepackage[]{microtype}
  \UseMicrotypeSet[protrusion]{basicmath} % disable protrusion for tt fonts
}{}
\makeatletter
\@ifundefined{KOMAClassName}{% if non-KOMA class
  \IfFileExists{parskip.sty}{%
    \usepackage{parskip}
  }{% else
    \setlength{\parindent}{0pt}
    \setlength{\parskip}{6pt plus 2pt minus 1pt}}
}{% if KOMA class
  \KOMAoptions{parskip=half}}
\makeatother
\usepackage{longtable,booktabs,array}
\usepackage{calc} % for calculating minipage widths
% Correct order of tables after \paragraph or \subparagraph
\usepackage{etoolbox}
\makeatletter
\patchcmd\longtable{\par}{\if@noskipsec\mbox{}\fi\par}{}{}
\makeatother
% Allow footnotes in longtable head/foot
\IfFileExists{footnotehyper.sty}{\usepackage{footnotehyper}}{\usepackage{footnote}}
\makesavenoteenv{longtable}
\setlength{\emergencystretch}{3em} % prevent overfull lines
\providecommand{\tightlist}{%
  \setlength{\itemsep}{0pt}\setlength{\parskip}{0pt}}
% paper_preamble.tex — shared LaTeX preamble for the PCMA/SI paper set.
%
% Compiled with lualatex (NOT pdflatex) because several papers contain
% Bengali script: the CLL paper realises one cognitive state into
% English, Bengali and Banglish, so the Bengali output is EVIDENCE and
% cannot be transliterated away without destroying the point being made.
%
% Font selection is NOT done here. pandoc >= 3 has native
% `mainfontfallback` / `monofontfallback` variables that emit the
% luaotfload fallback registration in the correct position relative to
% unicode-math; doing it by hand in this header lands after unicode-math
% and breaks font loading outright. See tools/build_papers.sh.

\usepackage{amsmath}
\usepackage{amssymb}
\usepackage{mathtools}   % \xleftrightarrow etc. used in cll.md §5
\usepackage{booktabs}

% ---------------------------------------------------------------------
% Unicode symbols appearing in PROSE rather than inside $...$.
%
% The papers were written as markdown and read as markdown, so they use
% real characters — "the ω-boost", "α-child", "ΔR coherence" — instead of
% LaTeX macros. Under unicode-math those codepoints are math-mode-only
% and abort the build ("\symbballowed only in math mode"). Mapping each
% to \ensuremath keeps the markdown readable on GitHub AND compilable,
% which is the whole point of keeping one source.
% ---------------------------------------------------------------------
\usepackage{newunicodechar}

% Greek, lower case
\newunicodechar{α}{\ensuremath{\alpha}}
\newunicodechar{β}{\ensuremath{\beta}}
\newunicodechar{γ}{\ensuremath{\gamma}}
\newunicodechar{δ}{\ensuremath{\delta}}
\newunicodechar{ε}{\ensuremath{\varepsilon}}
\newunicodechar{η}{\ensuremath{\eta}}
\newunicodechar{θ}{\ensuremath{\theta}}
\newunicodechar{κ}{\ensuremath{\kappa}}
\newunicodechar{λ}{\ensuremath{\lambda}}
\newunicodechar{μ}{\ensuremath{\mu}}
\newunicodechar{π}{\ensuremath{\pi}}
\newunicodechar{ρ}{\ensuremath{\rho}}
\newunicodechar{σ}{\ensuremath{\sigma}}
\newunicodechar{τ}{\ensuremath{\tau}}
\newunicodechar{φ}{\ensuremath{\varphi}}
\newunicodechar{χ}{\ensuremath{\chi}}
\newunicodechar{ψ}{\ensuremath{\psi}}
\newunicodechar{ω}{\ensuremath{\omega}}
\newunicodechar{ν}{\ensuremath{\nu}}
\newunicodechar{ξ}{\ensuremath{\xi}}
\newunicodechar{ζ}{\ensuremath{\zeta}}
\newunicodechar{ι}{\ensuremath{\iota}}

% Greek, upper case
\newunicodechar{Δ}{\ensuremath{\Delta}}
\newunicodechar{Φ}{\ensuremath{\Phi}}
\newunicodechar{Σ}{\ensuremath{\Sigma}}
\newunicodechar{Ψ}{\ensuremath{\Psi}}
\newunicodechar{Ω}{\ensuremath{\Omega}}
\newunicodechar{Γ}{\ensuremath{\Gamma}}
\newunicodechar{Λ}{\ensuremath{\Lambda}}
\newunicodechar{Θ}{\ensuremath{\Theta}}
\newunicodechar{Π}{\ensuremath{\Pi}}

% Operators and relations
\newunicodechar{×}{\ensuremath{\times}}
\newunicodechar{·}{\ensuremath{\cdot}}
\newunicodechar{±}{\ensuremath{\pm}}
\newunicodechar{−}{\ensuremath{-}}
\newunicodechar{≈}{\ensuremath{\approx}}
\newunicodechar{≤}{\ensuremath{\leq}}
\newunicodechar{≥}{\ensuremath{\geq}}
\newunicodechar{≠}{\ensuremath{\neq}}
\newunicodechar{∈}{\ensuremath{\in}}
\newunicodechar{∂}{\ensuremath{\partial}}
\newunicodechar{∇}{\ensuremath{\nabla}}
\newunicodechar{√}{\ensuremath{\sqrt{}}}
\newunicodechar{∞}{\ensuremath{\infty}}
\newunicodechar{∝}{\ensuremath{\propto}}
\newunicodechar{∑}{\ensuremath{\sum}}
\newunicodechar{∫}{\ensuremath{\int}}
\newunicodechar{ℝ}{\ensuremath{\mathbb{R}}}
\newunicodechar{ℕ}{\ensuremath{\mathbb{N}}}
\newunicodechar{‖}{\ensuremath{\|}}
\newunicodechar{⊕}{\ensuremath{\oplus}}
\newunicodechar{⊗}{\ensuremath{\otimes}}
\newunicodechar{∀}{\ensuremath{\forall}}
\newunicodechar{∃}{\ensuremath{\exists}}
\newunicodechar{¬}{\ensuremath{\neg}}
\newunicodechar{∧}{\ensuremath{\wedge}}
\newunicodechar{∨}{\ensuremath{\vee}}

% Arrows
\newunicodechar{→}{\ensuremath{\rightarrow}}
\newunicodechar{←}{\ensuremath{\leftarrow}}
\newunicodechar{↔}{\ensuremath{\leftrightarrow}}
\newunicodechar{↑}{\ensuremath{\uparrow}}
\newunicodechar{↓}{\ensuremath{\downarrow}}
\newunicodechar{⇒}{\ensuremath{\Rightarrow}}
\newunicodechar{⇐}{\ensuremath{\Leftarrow}}
\newunicodechar{⇔}{\ensuremath{\Leftrightarrow}}

% Super/subscripts written as literal characters
\newunicodechar{²}{\ensuremath{^{2}}}
\newunicodechar{³}{\ensuremath{^{3}}}
\newunicodechar{⁴}{\ensuremath{^{4}}}
\newunicodechar{⁵}{\ensuremath{^{5}}}
\newunicodechar{⁶}{\ensuremath{^{6}}}
\newunicodechar{₀}{\ensuremath{_{0}}}
\newunicodechar{₁}{\ensuremath{_{1}}}
\newunicodechar{₂}{\ensuremath{_{2}}}
\newunicodechar{₃}{\ensuremath{_{3}}}
\newunicodechar{ᵀ}{\ensuremath{^{\mathsf{T}}}}

% Status glyphs. The papers use emoji as status markers in tables
% (implemented / failed / warning). Mapped to typographic equivalents:
% a colour emoji font in lualatex is unreliable, and the MEANING is what
% carries — a reader needs to see "this one failed", not a red cross.
\newunicodechar{✅}{\ensuremath{\checkmark}}
\newunicodechar{✓}{\ensuremath{\checkmark}}
\newunicodechar{❌}{\ensuremath{\times}}
\newunicodechar{✗}{\ensuremath{\times}}
\newunicodechar{⚠}{\textbf{!}}
\newunicodechar{🧩}{\textbf{[piece]}}
\newunicodechar{📋}{\textbf{[list]}}
\newunicodechar{🟡}{\textbf{[partial]}}
\newunicodechar{🔴}{\textbf{[dead]}}
\newunicodechar{🟢}{\textbf{[live]}}
\newunicodechar{▼}{\ensuremath{\blacktriangledown}}
\newunicodechar{►}{\ensuremath{\blacktriangleright}}
\newunicodechar{□}{\ensuremath{\square}}
\newunicodechar{█}{\rule{0.6em}{0.9ex}}
\newunicodechar{•}{\textbullet}
\newunicodechar{†}{\textdagger}
\newunicodechar{‑}{-}
\newunicodechar{‐}{-}
\newunicodechar{…}{\ldots}
\newunicodechar{।}{\textbar}   % danda: rendered by the Bengali fallback in text
\newunicodechar{⁻}{\ensuremath{^{-}}}

% Zero-width formatting codepoints. These are invisible in an editor but
% are real characters to LaTeX: an emoji written as "⚠️" is TWO
% codepoints (U+26A0 U+FE0F), so declaring only the visible one leaves a
% variation selector with no definition and the build dies on a symbol
% nobody can see. Both are mapped to nothing.
\newunicodechar{^^^^fe0f}{}    % VARIATION SELECTOR-16
\newunicodechar{^^^^200d}{}    % ZERO WIDTH JOINER
\newunicodechar{^^^^200b}{}    % ZERO WIDTH SPACE

\usepackage{longtable}
\usepackage{array}

% Comparison matrices and cost ledgers are wide. Let them shrink to the
% text block rather than overflow the margin.
\usepackage{adjustbox}

% Code blocks: wrap long lines instead of running off the page. Several
% papers quote real C++ with long identifiers.
\usepackage{fvextra}
\fvset{breaklines=true,breakanywhere=true,fontsize=\small}

\setlength{\parskip}{0.5em}
\setlength{\parindent}{0pt}

\usepackage{fancyhdr}
\pagestyle{fancy}
\fancyhf{}
\fancyfoot[C]{\thepage}
\fancyhead[L]{\small\itshape Synthetic Intelligence — research lineage}
\fancyhead[R]{\small\itshape irfan@furylogic.com}
\renewcommand{\headrulewidth}{0.2pt}
\setlength{\headheight}{14pt}
\setlength{\headsep}{18pt}
\usepackage{bookmark}
\IfFileExists{xurl.sty}{\usepackage{xurl}}{} % add URL line breaks if available
\urlstyle{same}
\hypersetup{
  pdftitle={HNCA: Hierarchical Neural Constraint Architecture --- Cognition as Constraint Satisfaction across Two Coupled Scales},
  pdfauthor={Irfan Mahir --- Independent Researcher, Dhaka, Bangladesh --- irfan@furylogic.com},
  colorlinks=true,
  linkcolor={black},
  filecolor={Maroon},
  citecolor={Blue},
  urlcolor={[HTML]{0B5394}},
  pdfcreator={LaTeX via pandoc}}

\title{HNCA: Hierarchical Neural Constraint Architecture --- Cognition
as Constraint Satisfaction across Two Coupled Scales}
\author{Irfan Mahir --- Independent Researcher, Dhaka, Bangladesh ---
irfan@furylogic.com}
\date{}

\begin{document}
\maketitle

{
\hypersetup{linkcolor=}
\setcounter{tocdepth}{3}
\tableofcontents
}
\section{HNCA: Hierarchical Neural Constraint Architecture --- Cognition
as Constraint Satisfaction across Two Coupled
Scales}\label{hnca-hierarchical-neural-constraint-architecture-cognition-as-constraint-satisfaction-across-two-coupled-scales}

\textbf{Irfan Mahir} \emph{Independent Researcher, Dhaka, Bangladesh}
irfan@furylogic.com

\begin{quote}
\emph{``The engine is not searching for an answer. It is integrating
differential equations, and the answer is where the energy stops.''}
\end{quote}

\begin{center}\rule{0.5\linewidth}{0.5pt}\end{center}

\subsection{Abstract}\label{abstract}

We specify a two-scale substrate for cognition in which computation is
constraint satisfaction rather than search or prediction. The
\textbf{micro-layer} is a continuous-state cellular grid whose cells
carry physical variables and whose edges are differential equations
rather than learned weights. The \textbf{macro-layer} wraps clusters of
those cells into regions with distinct natural frequencies, coupled by
the Kuramoto model, reflecting the observation that brain regions
operate at different characteristic rhythms rather than as one
homogeneous field.

Processing has four stages: a perturbation is injected as a continuous
stimulus; the injected energy breaks the grid's equilibrium; the engine
integrates the constraint equations until the field settles; and the
resulting topological pattern \textbf{is} the thought, not a code for
one. No token is involved at any stage.

The macro-layer of this proposal is implemented and load-bearing in the
current engine. The micro-layer is not, and §6 states precisely which
part survived and why.

\begin{center}\rule{0.5\linewidth}{0.5pt}\end{center}

\subsection{1. Position}\label{position}

Two architectural options are commonly assumed to exhaust the space: a
flat cellular grid, or a Transformer over lexical tokens. Neither is
adequate for constraint-based cognition.

A \textbf{flat grid} has no notion of region. Every cell obeys the same
dynamics at the same characteristic rate, so there is no mechanism by
which one part of the system can operate on a different timescale from
another --- and differential timescale is what allows slow context to
modulate fast processing.

A \textbf{token Transformer} operates on discrete lexical units and
computes a distribution over continuations. Its geometry is a
fully-connected attention matrix, which is to say it has no geometry:
every unit is equidistant from every other, and spatial locality carries
no meaning.

What is required instead is a \textbf{physics engine for cognition} ---
a system with local interaction rules, spatial structure, and multiple
characteristic frequencies.

\begin{center}\rule{0.5\linewidth}{0.5pt}\end{center}

\subsection{2. The Micro-Layer: Continuous-State Cellular
Dynamics}\label{the-micro-layer-continuous-state-cellular-dynamics}

The foundational layer is a multi-dimensional continuous grid. A cell
does not contain a string or an embedding; it represents a micro-column
and carries a state vector of physical quantities: membrane potential
\(V_m\), a recovery variable \(u\), and a local field potential.

The edges are \textbf{differential equations}, not learned parameters.
The cell update implements a spiking model natively --- for instance the
Izhikevich formulation:

\[
\frac{dV}{dt} = 0.04V^2 + 5V + 140 - u + I
\]

\[
\frac{du}{dt} = a(bV - u)
\]

The significance is what is \emph{absent}. There is no embedding table,
no attention, and no parameter fitted by gradient descent. The
constraints are stated, not learned, and the dynamics they induce are
the computation.

\begin{center}\rule{0.5\linewidth}{0.5pt}\end{center}

\subsection{3. The Macro-Layer: Hierarchical
Oscillators}\label{the-macro-layer-hierarchical-oscillators}

A brain is not a homogeneous grid. Cortex, thalamus and amygdala operate
at different characteristic frequencies, and that heterogeneity is
functional rather than incidental.

Clusters of micro-cells are therefore wrapped into \textbf{macro-nodes},
each with a natural frequency \(\omega_i\), and coupled through the
Kuramoto model. The engine continuously solves for phase across the
network:

\[
\frac{d\theta_i}{dt} \;=\; \omega_i \;+\; \frac{K}{N}\sum_{j=1}^{N} \sin(\theta_j - \theta_i)
\]

Here \(\omega_i\) is the natural frequency of a region and \(K\) is the
coupling strength --- the analogue of the white-matter tracts connecting
regions.

Two properties follow. Regions with similar \(\omega\) \textbf{entrain}
easily and exchange influence readily. Regions with dissimilar
\(\omega\) resist entrainment, so a frequency separation acts as a
\emph{soft domain boundary} requiring no gating logic and no explicit
routing table. Domain isolation becomes a physical property rather than
a control structure.

\begin{center}\rule{0.5\linewidth}{0.5pt}\end{center}

\subsection{4. The Constraint Engine}\label{the-constraint-engine}

Input and output are both physical events.

\textbf{Stage 1 --- Inject perturbation.} Rather than passing text to an
encoder, a continuous stimulus (a current \(I\)) is injected into
specific sensory cells.

\textbf{Stage 2 --- Disequilibrium.} The injected energy breaks the
grid's equilibrium. Kuramoto oscillators fall out of phase and cell
states spike. The system is \emph{forced} to react; nothing decides to
respond.

\textbf{Stage 3 --- Continuous constraint solving.} The integration loop
runs. The engine is not searching a space of candidate answers --- it is
integrating differential equations across the edges. Energy propagates
bound strictly by the local rules.

\textbf{Stage 4 --- Attractor convergence.} The field settles into a new
stable configuration. \textbf{This topological pattern is the thought.}
A downstream decoder translates the settled state into text or action,
and that translation is a separate, explicitly-named cost rather than
part of the cognition.

\begin{center}\rule{0.5\linewidth}{0.5pt}\end{center}

\subsection{5. Comparison}\label{comparison}

\begin{longtable}[]{@{}
  >{\raggedright\arraybackslash}p{(\linewidth - 4\tabcolsep) * \real{0.3333}}
  >{\raggedright\arraybackslash}p{(\linewidth - 4\tabcolsep) * \real{0.3333}}
  >{\raggedright\arraybackslash}p{(\linewidth - 4\tabcolsep) * \real{0.3333}}@{}}
\toprule\noalign{}
\begin{minipage}[b]{\linewidth}\raggedright
Property
\end{minipage} & \begin{minipage}[b]{\linewidth}\raggedright
Token Transformer
\end{minipage} & \begin{minipage}[b]{\linewidth}\raggedright
Constraint engine
\end{minipage} \\
\midrule\noalign{}
\endhead
\bottomrule\noalign{}
\endlastfoot
Fundamental unit & lexical tokens & continuous state variables \\
Mechanics & probabilistic next-symbol prediction & deterministic
differential equations \\
Network geometry & fully-connected attention & spatially bounded
topological grid \\
Output generation & sampling from a distribution & relaxation into an
attractor \\
Timescales & one & many, set by \(\omega_i\) per region \\
\end{longtable}

The intended consequence is a change in what is being programmed. One
does not specify the logic; one specifies the \textbf{laws} the
substrate obeys, and the logic is whatever those laws force.

\begin{center}\rule{0.5\linewidth}{0.5pt}\end{center}

\subsection{6. What Was Implemented, and What Was
Not}\label{what-was-implemented-and-what-was-not}

This section is the paper's accounting against the running system, and
it corrects an earlier presentation in which the whole proposal was
described as live.

\textbf{The macro-layer is implemented and load-bearing.} Macro-nodes
with per-region natural frequencies, coupled by the Kuramoto update, are
in the current substrate (\texttt{struct\ MacroNode},
\texttt{physics.hpp}), and the frequency-separation-as-domain-boundary
property of §3 carries through to the sub-band ladder of the VGHT layer.
The claim that dissimilar \(\omega\) resists entrainment and thereby
isolates domains was tested and holds.

\textbf{The micro-layer is not implemented.} The current substrate does
not run Izhikevich cells on a continuous grid. Its Layer-1 unit is a
\textbf{mass-spring cavity} --- a damped oscillator coupled through
tunnels whose stiffness and damping are fixed by relation type. This is
a different micro-model presenting the same macro-interface, which is
precisely why the macro-layer transferred without modification.

\textbf{Why the substitution happened.} The spiking micro-layer was
specified for biological fidelity, and biological fidelity turned out
not to be load-bearing for the structural reasoning the engine is asked
to do. The mass-spring cavity is simpler, has an analytic impedance, and
admits a material table letting relation \emph{type} set transmission
properties. Nothing in §3 or §4 depended on the micro-model being
spiking.

The general lesson is worth stating: \textbf{the layered specification
survived a total replacement of its bottom layer.} That is evidence the
interface between the two scales was drawn in the right place, and it is
the most useful thing this paper produced.

\begin{center}\rule{0.5\linewidth}{0.5pt}\end{center}

\subsection{7. Limitations}\label{limitations}

\textbf{Attractor convergence is assumed, not guaranteed.} Section 4
asserts that the field settles. A driven nonlinear system need not
converge; it may oscillate indefinitely, or lock into a ground state and
refuse to leave. Ground-state locking was subsequently observed as a
real failure mode and required a dedicated escape mechanism.

\textbf{The decoder is unspecified.} Stage 4 defers the conversion of a
settled topology into text to ``a downstream decoder'', which is exactly
where the irreducible translation cost of the architecture lives.
Deferring it is not solving it, and naming it as deferred is the minimum
honest treatment.

\textbf{No account of acquisition.} Constraints are stated. How a system
acquires new constraints from experience lies outside this proposal and
was addressed separately by local Hebbian consolidation.

\begin{center}\rule{0.5\linewidth}{0.5pt}\end{center}

\subsection{8. Conclusion}\label{conclusion}

HNCA's contribution is the two-scale decomposition, and the specific
claim that differential natural frequency suffices for domain separation
without any routing logic. That claim survived into a substrate whose
micro-layer bears no resemblance to the one proposed here --- which is
the strongest available evidence that the decomposition, rather than
either individual layer, was the durable part.

\begin{center}\rule{0.5\linewidth}{0.5pt}\end{center}

\subsection{References}\label{references}

\begin{itemize}
\tightlist
\item
  Izhikevich, E. M. \emph{Simple Model of Spiking Neurons}, IEEE
  Transactions on Neural Networks, 2003.
\item
  Kuramoto, Y. \emph{Chemical Oscillations, Waves, and Turbulence},
  1984.
\item
  Mahir, I. \emph{From Artificial to Synthetic Intelligence.} §5.4c ---
  the substrate in which the macro-layer runs.
\item
  Mahir, I. \emph{Voltage-Gated Harmonic Topology (VGHT).} --- develops
  §3's frequency-separation argument into an explicit sub-band ladder.
\item
  Implementation: \texttt{src/physics.hpp} (\texttt{struct\ MacroNode},
  Kuramoto pass).
\end{itemize}

\end{document}
